\section{Positivity of the reciprocal of the diagonal Pad\'e denominator of the exponential}\label{sec:pade}

\subsection{The question}
For integers $k\ge0$, $m\ge1$ let
\[
q_{k,m}(z)=\sum_{j=0}^{m}(-1)^j\binom mj\frac{(k+m-j)!}{(k+m)!}\,z^j
\]
be the denominator of the $[k/m]$ Pad\'e approximant to $e^z$ at $0$ (normalised by $q_{k,m}(0)=1$). Fasi and Higham~\cite{FH} conjectured the following (their Conjecture 3.2) in connection with a forward-error bound for the matrix exponential.

\begin{conjecture}[Fasi--Higham~\cite{FH}]\label{conj:fh}
Let $k,m\in\N$. If $m$ is odd, then all the coefficients of the series expansion of $q_{k,m}(z)^{-1}$ at $0$ are positive.
\end{conjecture}

Scardino~\cite{Scardino} proved full positivity for $k=0$ and \emph{eventual} positivity (positivity of $c_n$ for all large $n$) whenever the positive real zero of $q_{k,m}$ is the unique zero of minimal modulus, a condition he shows holds for each fixed odd $m$ and all sufficiently large $k$; he also showed that the conjecture is false in general: with $c_n(k,m)$ the coefficients of $q_{k,m}^{-1}$, he gives exact counterexamples $c_{38}(66,5)<0$, $c_{37}(44,7)<0$, $c_{42}(45,9)<0$. He writes that ``the counterexamples are not diagonal, and the diagonal case $k=m$, which is the case relevant in practice, remains open.'' Here $q_{m,m}$ is a normalised reverse Bessel polynomial. We contribute two things: a computer-assisted proof of the diagonal case for $3\le m\le361$ (Theorem~\ref{thm:pade}), and a numerical description of where positivity first fails for fixed $m$ (Section~\ref{sec:kstar}). The general diagonal statement remains open.

\begin{theorem}[computer-assisted]\label{thm:pade}
For every odd $m$ with $3\le m\le 361$, all Taylor coefficients at $0$ of $1/q_{m,m}(z)$ are positive. (For $m=1$, $1/q_{1,1}(z)=\sum_n (z/2)^n$.)
\end{theorem}

\subsection{The certificate}
Fix an odd $m$ and write $q=q_{m,m}$, $c_n$ for the coefficients of $1/q$.

\begin{lemma}\label{lem:cert}
Suppose $q$ has $m$ distinct zeros $\rho_1,\dots,\rho_m$, exactly one of which, $R=\rho_1$, is real and positive, and $|\rho_i|>R$ for $i\ge2$. Put $B_i=-1/q'(\rho_i)$, $S=\sum_{i\ge2}|B_i|$, $\rho_{\min}=\min_{i\ge2}|\rho_i|$. If $B_1>0$ and $N_0=\max\{0,\lceil\ln(S/B_1)/\ln(\rho_{\min}/R)\rceil\}$, then $c_n>0$ for all $n\ge N_0$. Consequently $c_n>0$ for all $n\ge0$ as soon as $c_0,\dots,c_{N_0-1}>0$.
\end{lemma}
\begin{proof}
Since $q(0)=1$ and $q$ has $m$ simple zeros, partial fractions give $1/q(z)=\sum_i \frac{1}{q'(\rho_i)(z-\rho_i)}$ (there is no polynomial part, the numerator having degree $0<m$), hence $c_n=\sum_iB_i\rho_i^{-n-1}$ for all $n\ge0$. For $n\ge N_0$ we have $n+1>\ln(S/B_1)/\ln(\rho_{\min}/R)$, i.e.\ $(\rho_{\min}/R)^{n+1}>S/B_1$, i.e.\ $S\rho_{\min}^{-n-1}<B_1R^{-n-1}$. Therefore
\[
c_n\ \ge\ B_1R^{-n-1}-\sum_{i\ge2}|B_i||\rho_i|^{-n-1}\ \ge\ B_1R^{-n-1}-S\rho_{\min}^{-n-1}>0 .\qedhere
\]
\end{proof}
(In the computation we take $N_0$ one larger than the displayed value to absorb rounding; this only makes the check stronger.)

\subsection{The computation}
For each odd $m$ in the range we carried out, in rigorous ball arithmetic (Arb through python-flint~\cite{Arb}):
\begin{enumerate}[leftmargin=2em]
\item isolation of the $m$ complex zeros of $q$ in disjoint balls; exactly one ball is real and positive, giving $R$; for each of the other balls a certified lower bound of the modulus exceeding $R$;
\item the residues $B_i=-1/q'(\rho_i)$ (certified balls), with $B_1>0$ certified, and the resulting bound $N_0$ of Lemma~\ref{lem:cert};
\item positivity of $c_0,\dots,c_{N_0-1}$ by inverting the power series $q$ in ball arithmetic, with the working precision raised until every one of these coefficients is a ball strictly above $0$.
\end{enumerate}
The working precision needed grows with $m$ (from 512 bits for $m\le11$ to 8192 bits for $m>300$), because the ratio $\rho_{\min}/R$ tends to $1$: it is $1.094$ for $m=3$, $1.008$ for $m=11$ and $1+3.5\cdot10^{-5}$ at $m\approx 240$ (Table~\ref{tab:pade}). The runs for $3\le m\le 361$ were recorded in log files (one line per $m$) in the repository; every one of the $180$ odd values of $m$ in this range appears with the result ``proved'', and no value is missing or failed. The largest threshold is $N_0=372{,}888$ at $m=361$ (at that point $c_n$ for $n<N_0$ are verified by series inversion with 8192 bits).

In addition, a second, independent exact certification of the first two items in the list was run: for every odd $3\le m\le361$, $\gcd(Q,Q')=1$ for the integer polynomial $Q=(2m)!\,q$ (so the zeros are simple), the zeros were isolated by Arb's certified routine for integer polynomials (which also certifies the real ones), exactly one zero is real and positive, and all other zeros have certified modulus larger than $R$.

\begin{table}[h]
\centering
\caption{Data from the certification runs (selected odd $m$). $R$ is the positive real zero of $q_{m,m}$, $\rho_{\min}/R$ the gap to the next zero in modulus, $N_0$ the threshold of Lemma~\ref{lem:cert}.}\label{tab:pade}
\begin{tabular}{rrrr}
\toprule
$m$ & $R$ & $\min_{i\ge2}|\rho_i|/R$ & $N_0$\\
\midrule
3 & 4.64437 & 1.09446 & 2\\
11 & 15.2447 & 1.00798 & 128\\
101 & 134.537 & 1.000102 & 22\,697\\
201 & 267.086 & 1.0000258 & 103\,850\\
301 & -- & -- & 251\,082\\
361 & -- & -- & 372\,888\\
\bottomrule
\end{tabular}
\end{table}

\paragraph{What is and is not proved.} Theorem~\ref{thm:pade} is a theorem for each listed $m$ individually; the proof consists of Lemma~\ref{lem:cert} together with the certified computations above, which are rigorous in the sense of ball arithmetic but are performed by a program (in the repository, \texttt{results/pade\_positivity}). It does not give the statement for all $m$. Moreover the gap $\rho_{\min}/R\to1$ and the growth of $N_0$ (the ratio $N_0/m^2$ is $2.2,\,2.6,\,2.8,\,2.9$ for $m=101,201,301,361$) indicate that this method cannot be pushed to all odd $m$; a proof for all $m$ needs a different argument.

\subsection{Where positivity first fails}\label{sec:kstar}
For fixed odd $m$ let $k^*(m)$ be the least $k$ for which a negative coefficient appears. We computed $k^*(m)$ by bisection in $k$ with interval arithmetic (each $c_n$, $n\le5k+300$, certified to have a definite sign), assuming that the set of $k$ for which negative coefficients occur is an interval $[k^*,\infty)$; we verified $k^*-1$ has no negative coefficient in that range and that $k^*,\dots,k^*+19$ all do, and checked contiguity directly only for $15\le m\le71$ and $k<400$. For small $m$ the values agree with exact rational arithmetic: $k^*(5)=66$ (first negative coefficient $n=38$), $k^*(7)=44$ ($n=37$), $k^*(9)=42$ ($n=41$), $k^*(19)=54$ ($n=68$); the first two are Scardino's examples, and $(k,m,n)=(42,9,41)$ is a smaller counterexample than his $(45,9,42)$. For $m=3$ we found no negative coefficient for any $k<1500$ (checking $n\le4k+200$). For larger $m$ the computation gives
\[
k^*(m)-2m\in\{12,13,14\}\qquad(109\le m\le 1009\ \text{sampled}),
\]
e.g.\ $k^*(109)=230$, $k^*(209)=430$, $k^*(509)=1031$, $k^*(1009)=2032$, and the first negative coefficient occurs near $n\approx k+m-3$. For $m=109$ and $m=209$ the thresholds were re-checked with exact rational arithmetic around $k^*$ (no negative coefficient among $c_n$, $n\le k+m+10$ for $k=229,429$; negative coefficients $c_{337}(230,109)$ and $c_{637}(430,209)$). In particular the diagonal $k=m$ is far from the boundary of the failure region for large $m$ ($k^*/m\to2$ numerically). This law is a numerical observation, not a theorem.

\begin{question}
Is $1/q_{m,m}$ positive for every odd $m$? Is $k^*(m)=2m+O(1)$?
\end{question}
